\subsection{The cover coefficient is the actual finite Fourier transform}
\label{subsec:ftb-cover-finite-fourier}

This subsection proves the coefficient identity used immediately
before the root projection. It keeps the algebraic cover and the
markings of the source. In particular its phase is determined by
the two actual roots of unity and by the two dual markings.

Write \(d\) for the supplied level-cover datum and \(p=d.pd>0\).
The two levels have integers \(M,M_2\), with
\(M_2=pM\) and \(\gcd(p,M)=1\). Put
\[
 P_0=\operatorname{Mfr}(l),\quad Q_0=l.\operatorname{Mb},
 \quad I=d.I,\quad J=\bar I,\quad
 \mathfrak m=\operatorname{Mfr}(l_2)=P_0J,\quad
 Q_2=l_2.\operatorname{Mb}=Q_0I.
\]
The cover and \(\mathrm{LevTie}\) give
\[
 I+J=\mathcal O_E,\quad IJ=(p),\quad
 P_0+J=\mathcal O_E,\quad
 \mathfrak m+\bar{\mathfrak m}=\mathcal O_E,\quad
 (2)+\mathfrak m=\mathcal O_E.                         \tag{CFo1}
\]
The rank is the literal \(r=Dw.p+Dw.q\), and \(r\leq n\).
Let \(D\) be the actual Hermitian level datum. Thus it has
\(\beta_i\in X.\Lambda\), its subgroup \(D.J\) and character
\(D.\lambda\), and \(\mathrm{LevTie}\) identifies
\(D.\mathfrak m=\mathfrak m\) and \(D.w=H.w_0\).

We use the original localized congruence
\[
 u\equiv_{\mathfrak a}v
 \quad\Longleftrightarrow\quad
 \exists t\in\mathcal O_E,\ (t)+\mathfrak a=\mathcal O_E,\
                          t(u-v)\in\mathfrak a X.\Lambda .
                                                               \tag{CFo2}
\]
For \(L\in D.\mathrm{lats}\), there is a common denominator,
prime to \(\mathfrak m\bar{\mathfrak m}\), carrying \(L\) into
\(X.\Lambda\) and \(X.\Lambda\) into \(L\). Consequently, if
\(\mathfrak m\subset\mathfrak a\) and \(u-v\in L\), then
\[
                  u\equiv_{\mathfrak a}v
                   \quad\Longleftrightarrow\quad
                             u-v\in\mathfrak a L.         \tag{CFo3}
\]
Indeed multiply the localized equality by the common denominator
to put it in \(\mathfrak a L\), and then use a Bezout combination
of that denominator times \(t\) and \(\mathfrak a\) to cancel it.
The reverse implication is the same denominator clearing.
This is why the same \(L\in D.\mathrm{lats}\) is required.

Define \(F=\mathrm{datF}(D,P_0,r\leq n)\) literally as follows.
Its first \(r\) rows must be congruent modulo \(\mathfrak m\)
to \(\sum_jk'_{ji}\beta_j\), for a lift \(k'\) of some \(k\in D.J\);
its remaining rows must be congruent to zero modulo \(P_0\).
On this support \(F\) is \(D.\lambda(k)\), with the choice
convention of the source; elsewhere it is zero.
The level tie says that the matrices in \(D.J\) are congruent to
the identity modulo \(P_0\). Hence
\[
 F(y)\ne0\Longrightarrow
 \begin{cases}
 y_i\equiv_{P_0}\beta_i,&i<r,\\
 y_i\equiv_{P_0}0,&i\geq r.
 \end{cases}                                             \tag{CFo4}
\]
Also rowwise congruence modulo \(\mathfrak m\) preserves \(F\):
the identical reduction witnesses can be used in its definition.

Let \(m,m_2:\mathcal L\to L\) be the two \(\mathrm{IsMark}\)
maps, with the source compatibility
\[
 a\in\mathfrak m\mathcal L
       \Longrightarrow m_2(a)=m(a)+Mu\quad(u\in L).        \tag{CFo5}
\]
Let \(\zeta,\zeta_2\) denote the images of the actual primitive
\(M\)-th and \(M_2\)-th roots under the compatible coefficient-field
embeddings. Let \(\delta:\mathcal L\to L\) be the old dual marking:
\[
 \zeta^{Dw.EF(x,c_0a)}
           =e^{-2\pi iEF_X(\delta x,M^{-1}m(a))}
                                      \quad(a\in P_0\mathcal L),
                                                               \tag{CFo6}
\]
where \(c_0\in P_0,\ c_0-1\in Q_0\) is the source \(cSep(l)\).
Retain the higher-level dual marking \(\delta_2:\mathcal L\to L\),
with
\[
 \zeta_2^{Dw.EF(x,c_2a)}
           =e^{-2\pi iEF_X(\delta_2x,M_2^{-1}m_2(a))}
                                     \quad(a\in\mathfrak m\mathcal L),
                                                               \tag{CFo7}
\]
where \(c_2\in\mathfrak m,\ c_2-1\in Q_2\). Its additional
properties are exactly
\[
 L\subset\delta_2(\mathcal L)+\mathfrak mL,\qquad
 \delta_2(a)\in\mathfrak mL
                        \Longleftrightarrow a\in\mathfrak m\mathcal L.
                                                               \tag{CFo8}
\]
These are the output of the actual
\srcid{dual_half_level_marking_exists_and_match},
not independent arbitrary dual vectors.

Let \(x^0\in\mathcal L^n\) be the source adapted base class:
\[
 p\delta(x^0_i)\equiv_{P_0}\beta_i\ (i<r),\qquad
                          x^0_i\in P_0\mathcal L\ (i\geq r).
                                                               \tag{CFo9}
\]
The coefficient function \(\mathrm{ph}\) on
\(\mathcal X=(\mathcal L/J\mathcal L)^n\) is part of the supplied
\(\mathrm{PhiPin}\), and \(\mathrm{LevTie}\) retains its actual
compatibility with \(F\): for a map
\(\delta_J:\mathcal L\to L\) surjective modulo \(J\), there is
one unit scalar \(t\), of absolute value one, such that
\[
 \sigma(\mathrm{ph}([\mathbf a]))
                      =tF(y)                            \tag{CFo10}
\]
whenever \(y_i\equiv_J\delta_J(a_i)\), \(y_i\equiv_{P_0}\beta_i\)
for \(i<r\), and \(y_i\equiv_{P_0}0\) for \(i\geq r\).
This explicitly states the common-model compatibility still
required of the supplied cover.

\begin{srlemma}[The actual auxiliary marking and coefficient representatives]
\label{lem:ftb-cover-coefficient-reps}
For \(\delta_J=M\delta_2\),
\[
 \delta_J(\mathcal L)\subset P_0L,\quad
 L\subset\delta_J(\mathcal L)+JL,\quad
 \delta_J(a)\in JL\Longleftrightarrow a\in J\mathcal L.     \tag{CFo11}
\]
Choose any representatives \(a(x_i)\) for \(x_i\in\mathcal L/J\mathcal L\),
and put
\[
                    z_{x,i}=p\delta(x^0_i)+\delta_J(a(x_i)).
                                                               \tag{CFo12}
\]
Then \(z_{x,i}\in L\), and the scalar \(t\) in (CFo10), chosen
for this \(\delta_J\), satisfies
\[
                          \sigma(\mathrm{ph}(x))=tF(z_x).
                                                               \tag{CFo13}
\]
The \(z_x\)'s are distinct modulo \(\mathfrak mL\) row by row.
Every \(w\) with rows in \(L\) and \(F(w)\ne0\) is congruent
row by row modulo \(\mathfrak mL\) to some \(z_x\).
\end{srlemma}

\begin{proof}
Since \(M\in P_0\), the first assertion of (CFo11) follows.
Choose integers \(a,b\) with \(ap+bM=1\).
By (CFo8), \(v=\delta_2(u)+w\), \(w\in\mathfrak mL\), so
\[
                 v=\delta_J(bu)+ap\,\delta_2(u)+w.
\]
Here \(p\in J\), giving the second assertion.
If \(\delta_J(u)\in JL\), it also lies in \(P_0L\), so lies
in \(P_0L\cap JL=P_0JL=\mathfrak mL\).
Apply (CFo8) to \(Mu\) and cancel \(M\) modulo \(J\), using
\(ap+bM=1\) and \(p\in J\). This gives \(u\in J\mathcal L\).
The converse follows from \(\delta_J(\mathcal L)\subset L\).

The vector \(p\delta(x^0_i)\) belongs to \(JL\).
The second summand of (CFo12) belongs to \(P_0L\).
Thus (CFo9) and (CFo3) verify all the conditions of (CFo10),
giving (CFo13). For the lower rows, use the actual
\(x^0_i\in P_0\mathcal L\), rather than a condition only on
their reductions.

If \(z_x-z_{x'}\) is rowwise in \(\mathfrak mL\subset JL\),
the last assertion of (CFo11) gives \(x=x'\).
Conversely, for a supported \(w\), choose \(a_i\) with
\(w_i-\delta_J(a_i)\in JL\), by (CFo11).
Since \(p\delta(x^0_i)\in JL\), the difference
\(w_i-z_{[a],i}\) is in \(JL\).
It is in \(P_0L\) by (CFo4), (CFo9), and (CFo3).
Coprimality \(P_0+J=\mathcal O_E\) now puts it in
\(P_0L\cap JL=\mathfrak mL\).
\end{proof}

\begin{srlemma}[The original tensor-quotient reindexing]
\label{lem:ftb-cover-tensor-reindex}
Let \(\mathfrak b+\mathfrak m=\mathcal O_E\). Suppose an integral
\(s\) satisfies \(s-1\in\mathfrak m\),
\(s\Lambda\subset\mathcal O_E^n\), and
\(s\mathcal O_E^n\subset\Lambda\).
Let \(G:\operatorname{Mat}_{n,r}(E)\to\mathbb C\) satisfy
\[
 \begin{aligned}
 &G(y)=G(sy) &&(y\in T(\Lambda,\mathfrak bL));\\
 &G(w)=G(w')&&
    (w_i\in L,\ w_i-w'_i\in\mathfrak mL\text{ for all }i).
 \end{aligned}                                             \tag{CFo14}
\]
Suppose \(z_x\) have integral rows, are distinct rowwise modulo
\(\mathfrak mL\), and meet every supported such congruence class.
Then
\[
 \sum_{c\in T(\Lambda,\mathfrak bL)/
                    T(\Lambda,\mathfrak b\mathfrak mL)}
                 G(c_{\mathrm{out}})=\sum_xG(z_x).       \tag{CFo15}
\]
\end{srlemma}

\begin{proof}
Choose \(e\in\mathfrak b,e'\in\mathfrak m\) with \(e+e'=1\).
A tensor-quotient class of \(y\) is characterized by the row
classes of \(sy\) modulo \(\mathfrak mL\).
The forward implication follows by clearing rows of the
tensor \(y-y'\).
For the reverse implication let
\(w=y-y'\in T(\Lambda,\mathfrak bL)\), with
\(sw_i\in\mathfrak mL\).
The same rows lie in \(\mathfrak bL\), and coprimality gives
\(sw_i\in\mathfrak b\mathfrak mL\). As \(se_i\in\Lambda\),
\[
                   s^2w\in T(\Lambda,\mathfrak b\mathfrak mL).
\]
Also \((s+1)(s-1)w\) is in that tensor group, since
\(s-1\in\mathfrak m\). Their difference is \(w\), proving
the characterization.

Assign \(x\) the class \(c_x=[sez_x]\).
It belongs to the required tensor group because \(se_i\in\Lambda\)
and \(ez_{x,i}\in\mathfrak bL\).
Moreover
\[
 s^2e-1=(s+1)e(s-1)-e'\in\mathfrak m,
\]
so the row class of \(s(sez_x)\) is that of \(z_x\).
The characterization therefore proves that \(x\mapsto c_x\)
is injective and that
\(c=c_x\) exactly when \(sc_{\rm out}-z_x\) is rowwise in
\(\mathfrak mL\).
By (CFo14), \(G(c_x)=G(z_x)\).
Every class with nonzero \(G(c_{\rm out})\) is one of the \(c_x\),
by the support hypothesis applied to \(sc_{\rm out}\).
Removing the zero terms gives (CFo15).
The argument is also the source's \(\mathrm{finsum}\) proof when
the indexing sets have not yet been equipped with finite
enumerations.
\end{proof}

\begin{srproposition}[The actual coefficient Fourier identity]
\label{prop:ftb-cover-coefficient-fourier}
For the supplied cover, markings, \(\mathrm{PhiPin}\) and
\(\mathrm{LevTie}\) above, there is a scalar \(t\) with \(|t|=1\),
chosen before \(\Lambda,N,Y,q\), such that
\[
 \begin{split}
 \sigma_2(d.w(x^0,q))={}&t
   \sum_{c\in T(\Lambda,2\bar{\mathfrak m}L)/
                    T(\Lambda,2\bar{\mathfrak m}\mathfrak mL)}
    F(c_{\rm out})
       e^{-2\pi i\sum_iEF_X(c_{{\rm out},i},Y_{q,i})}.
 \end{split}                                             \tag{CFo16}
\]
Here \(\Lambda\in\operatorname{Vset}(E,n,M_2)\),
\(N\) satisfies \(\mathrm{NOK}(\Lambda,N,M_2)\), and \(Y\)
satisfies the literal \(\mathrm{IsRoots}(l_2,L,m_2,N,Y)\).
There is no parity or large-rank restriction in this coefficient
identity beyond its displayed \(r\leq n\).
\end{srproposition}

\begin{proof}
Use the scalar \(t\) from (CFo13), independent of these later
choices. Put
\[
              \mathfrak b=2\bar{\mathfrak m},\qquad
 G(y)=F(y)\exp\{-2\pi i\sum_iEF_X(y_i,Y_{q,i})\}.
\]
Coprimality in (CFo1) gives
\(\mathfrak b+\mathfrak m=\mathcal O_E\).
Choose an integer \(s\equiv1\pmod {M_2}\) clearing \(\Lambda\)
in both directions: if \(h\) in (RP3) is prime to \(M_2\),
choose \(a,b\in\mathbb Z\) with \(ah+bM_2=1\) and take \(s=ah\).
Thus \(s-1\in\mathfrak m\), as required in
Lemma~\ref{lem:ftb-cover-tensor-reindex}.

We verify (CFo14), including the dependence on the actual roots.
First \(s(y_i-(sy)_i)=-(s-1)(sy)_i\in\mathfrak mL\),
so denominator clearing and (CFo2) make \(F(y)=F(sy)\).
For the phase, the root assumptions imply
\(Q_2Y_{q,i}\subset L\), hence \(M_2Y_{q,i}\in L\).
If \(a\in E\) has \(sa,Na\) integral, the
\(Q_2\)-power condition gives \(M_2^jY_{q,i}=2Nu\), \(u\in L\).
Choose integers \(A,B\) with \(As+BM_2^j=1\).
Then the exact identity
\[
 aM_2Y_{q,i}
       =A(sa)M_2Y_{q,i}+B(2M_2)(Na)u
\]
shows that \(aM_2Y_{q,i}\in L\).
For \(v\in\Lambda\), both \(sv_i,Nv_i\) are integral.
Apply this observation to \(a=\bar v_i\), move the coefficient
by the adjoint formula (RP8), and use selfduality of \(L\).
On a pure tensor \(v\otimes\ell\), and hence on
\(y\in T(\Lambda,\mathfrak bL)\), it follows that
\(EF_X(y_i,M_2Y_{q,i})\in\mathbb Z\).
Since \(s-1\) is an integer multiple of \(M_2\),
the phases of \(y\) and \(sy\) agree.

If \(w_i-w'_i\in\mathfrak mL\), congruence invariance
already gives \(F(w)=F(w')\).
For \(a\in\mathfrak m\), \(\bar a\in Q_2\), so
\(\bar aY_{q,i}\in L\).
The adjoint formula and selfduality give
\(EF_X(w_i-w'_i,Y_{q,i})\in\mathbb Z\).
This proves the other phase invariance in (CFo14).
Lemma~\ref{lem:ftb-cover-coefficient-reps} supplies the integral,
distinct and supported \(z_x\). Thus (CFo15) applies.

It remains to identify \(G(z_x)\) with the cover coefficient,
rather than assuming its phase. For this write
\(a_i\in\mathfrak m\mathcal L\) for the root representative
of \(q_i\); thus
\(Y_{q,i}-M_2^{-1}m_2(a_i)\in L\).
Let \(c_I\in J,\ c_I-1\in I\) be the actual \(cHalf(H,I)\).
Modulo integral pairings,
\[
 \begin{split}
 EF_X(z_{x,i},Y_{q,i})\equiv{}&
 EF_X(\delta(x^0_i),M^{-1}m(a_i))\\
 &+M\,EF_X(\delta_2(a(x_i)),M_2^{-1}m_2(a_i)).
 \end{split}                                             \tag{CFo17}
\]
Indeed discard the integral difference in \(Y\), use
\(M_2=pM\), expand (CFo12), and use (CFo5) to replace
\(M^{-1}m_2(a_i)\) by \(M^{-1}m(a_i)\).
Both discarded pairings have two arguments in \(L\).

The old and new dual identities (CFo6)--(CFo7) consequently give
\[
 e^{-2\pi iEF_X(z_{x,i},Y_{q,i})}
       =\zeta^{Dw.EF(x^0_i,c_0a_i)}
          (\zeta_2^M)^{Dw.EF(a(x_i),c_2a_i)}.             \tag{CFo18}
\]
Let \(\widetilde a_i=\operatorname{levRep}(q_i)\).
Since \(a_i-\widetilde a_i\in Q_2\mathcal L\subset Q_0\mathcal L\),
\(c_0(a_i-\widetilde a_i)\in M\mathcal L\); the first exponent
may therefore use \(\widetilde a_i\).
For the second one,
\[
                   c_2a_i-c_I\widetilde a_i
                                      \in p\mathcal L.   \tag{CFo19}
\]
To check this, \(c_2-c_I\) belongs to \(I\cap J=IJ=(p)\):
both elements are \(1\) modulo \(I\) and belong to \(J\).
Also \(c_I(a_i-\widetilde a_i)\in JI\mathcal L\).
Since \((\zeta_2^M)^p=1\), (CFo19) changes no phase.
Finally, the source roots of unity give exactly
\[
  \zeta_2^M
       =\sigma_2\!\left(\zeta(M'_2)^{M'_2/p}\right).       \tag{CFo20}
\]
This is an equality of the specified powers, obtained from
\(M_2=pM\) and \(M_2\mid M'_2\).

The last clause of \(\mathrm{PhiPin}\) is literally
\[
 \sigma_2(d.eF(d.\phi_\delta(q)))
   =\sum_{x\in\mathcal X}\sigma(\mathrm{ph}(x))
      (\zeta_2^M)^{
             \sum_iDw.EF(a(x_i),c_I\widetilde a_i)}.       \tag{CFo21}
\]
By definition \(d.w(x^0,q)\) multiplies this by
\(\zeta^{\sum_iDw.EF(x^0_i,c_0\widetilde a_i)}\).
Substituting (CFo13) and (CFo18)--(CFo20) therefore gives
\[
                    \sigma_2(d.w(x^0,q))=t\sum_xG(z_x).
\]
The quotient reindexing just proved converts this into (CFo16).
Its sets are finite in this application: \(\mathcal X\) is a
finite integral-lattice quotient; the tensor quotient is a
quotient of a finite free abelian group killed by the positive
integer \(M_2\).
\end{proof}

\paragraph{Exact frozen-source correspondence.}
The endpoint is
\srcid{GlobalTheory.Dict.CoverFourier.w\_eq\_sum}, also exposed as
\srcid{GlobalTheory.T.coverFourierCoeff\_eq\_finite\_fourier\_sum},
in
\srcid{src/Development/AbelianVarietyTheory/Proofs/SiegelThetaTransformation/CoverFourierCoeffEqFiniteFourierSum.lean}
(\srcid{74576b27f7c3dd7a54f9c4d66edcd5931e6c9cb4}),
at tag \srcid{alpoge-source-7cbf7a0c}.
The tensor reindexing is its \srcid{finsum\_tens};
(CFo11) is its \(\delta_J=M\delta_2\) construction;
the two denominator-phase arguments are its
\srcid{root\_mem}, \srcid{root\_mem'}, \srcid{EF\_int\_tens}
and \srcid{EF\_int\_smul}.
The exact integral comparison (CFo19) is its \srcid{pd\_dvd},
and (CFo20) is \srcid{zetaM\_pow}.
The cover coefficient (CFo21), compatible embeddings, and
\srcid{PhiPin} are defined in
\srcid{src/Formalism/LocalToGlobal/ProjectiveLattices/ThShort.lean}
(\srcid{ba33f03086ab6eb7e86ad7e31fe5d9484953af5e}).
The localized congruence, \(D.\mathrm{lats}\), and \(D.\mathrm{Red}\)
are in
\srcid{src/Development/AbelianVarietyTheory/Definitions/SpanningFamilies/IsSystem.lean}
(\srcid{ae90008828f2ba0ca764c4a5a36a0625faa93079});
the marking, root and dual identities are in its sibling
\srcid{IsStabForm.lean}
(\srcid{ca69affd0453770815104e7bfc7b4a93c7ef988e}).
The exact \srcid{LevTie} and adapted base class are in
\srcid{IsBaseClassLevelAdapted.lean}
(\srcid{24574b3c56538e14a28fa38ddc4e9b034c240057}).
The existence of the higher dual marking is the separate
\srcid{DualHalfLevelMarkingExistsAndMatch.lean}
(\srcid{ca7df4a980bb13ba448a4268b4ac12dd635d7c9f});
its output (CFo7)--(CFo8) has been stated explicitly here.
Thus this subsection proves the downstream coefficient identity
(RS8) from the actual cover and matching inputs. It does not
obtain those inputs by naming their interfaces.
